% Agent-Time: Data-Driven Task Duration Estimation for AI Agent Workflows
% Target venue: MSR 2027 (Mining Software Repositories) — Tools Track
% Paper ID: P23-AT
% Status: SEED — structure complete, \todo{} markers for measurement pass
%
% Thesis: AI agent task duration distributions differ structurally from human
% developer estimates. Training on AI agent production work logs (done-log +
% git history) produces better estimates than story points or naive mean.
% The distribution SHAPES are themselves a contribution: fat-tailed for
% research/investigation tasks, narrow for citation fixes and scaffolding.
%
% What's built:
%   D.030: git log + done-log \to task-type duration distributions
%   holo_estimate_task_duration: live MCP tool, operational
%   Data: HoloMesh done-log (all agent completions, timestamped, with board IDs)
%         + HoloScript git history (timestamps, file paths, LOC deltas)
%
% Novelty: First system trained on AI agent production logs, not human estimates.

\documentclass[sigconf]{acmart}
\usepackage{booktabs}
\usepackage{xcolor}
\usepackage{listings}
\usepackage{pgfplots}
\pgfplotsset{compat=1.18}

\newcommand{\todo}[1]{\textcolor{red}{\textbf{[TODO: #1]}}}
\newcommand{\holoscript}{HoloScript}
\newcommand{\holomesh}{HoloMesh}
\newcommand{\agenttime}{\textsc{AgentTime}}

\lstset{basicstyle=\scriptsize\ttfamily,breaklines=true,frame=single,
        numbers=left,numberstyle=\tiny\color{gray},xleftmargin=1.5em}

\title{\agenttime{}: Data-Driven Task Duration Estimation\\for AI Agent Workflows}

\author{Joseph Krzywoszyja}
\affiliation{\institution{HoloScript AI Lab}}
\email{josephtaxwise@gmail.com}

\begin{abstract}
Estimating task duration is essential for AI agent project management.
Human-derived approaches (story points, T-shirt sizing) assume human
cognitive load and working memory limits that do not apply to AI agents.
We introduce \agenttime{}, a system that mines AI agent production work
logs — the HoloMesh done-log and \holoscript{} git history — to learn
per-task-type duration distributions.
We find that AI agent task duration distributions differ structurally from
human estimates: research and investigation tasks follow heavy-tailed
distributions (occasional very long sessions are informative, not outliers),
while mechanical tasks (citation additions, scaffolding) follow narrow
distributions close to their median.
\agenttime{} is deployed as \texttt{holo\_estimate\_task\_duration}, a live
MCP tool integrated with the \holomesh{} board claim workflow.
\todo{Report: N completions across M task types, retrospective accuracy
vs. naive mean baseline.}
\end{abstract}

\keywords{AI agents, project management, task estimation, mining software repositories, multi-agent systems}

\begin{document}
\maketitle

% ─────────────────────────────────────────────────────────────────────────────
\section{Introduction}

AI agents executing software engineering tasks need duration estimates for
two reasons: (1) agents self-schedule across parallel workstreams and need
to know when to hand off vs. continue; (2) human oversight requires
knowing whether a task is running long (potential loop) or on track.

Human project management uses story points~\cite{cohn2005} and planning
poker, calibrated to human cognitive load, context-switching cost, and
meeting overhead.
None of these apply to AI agents.
An AI agent has no meeting overhead, has effectively infinite working memory
(up to context window), and its ``cognitive load'' is better modeled by
token budget and tool-call depth than by complexity.

We observe that AI agent task duration is better predicted by:
\begin{enumerate}[leftmargin=*,topsep=2pt,itemsep=1pt]
  \item \textbf{Task type}: research tasks systematically take longer than
        implementation tasks of the same stated complexity.
  \item \textbf{File scope}: tasks touching more files have heavier-tailed
        distributions.
  \item \textbf{Session history}: first-task-of-session takes longer
        (warm-up cost) than same-type tasks mid-session.
\end{enumerate}

\paragraph{Contributions.}
\begin{enumerate}[leftmargin=*,topsep=2pt,itemsep=1pt]
  \item \agenttime{}: a system mining done-log + git history to learn
        per-task-type duration distributions (Section~\ref{sec:system}).
  \item Empirical characterization of AI agent task duration distributions
        across \todo{M} task types (Section~\ref{sec:findings}).
  \item \texttt{holo\_estimate\_task\_duration}: deployed MCP tool with
        board integration (Section~\ref{sec:tool}).
  \item Comparison to naive mean baseline showing \agenttime{} achieves
        \todo{X\%} lower MAPE on held-out completions (Section~\ref{sec:eval}).
\end{enumerate}

% ─────────────────────────────────────────────────────────────────────────────
\section{Data Sources}

\subsection{HoloMesh Done-Log}

The \holomesh{} done-log is an append-only ledger of agent task completions.
Each entry contains:

\begin{itemize}[topsep=2pt,itemsep=1pt]
  \item \texttt{task\_id}: board task identifier
  \item \texttt{title}: free-text task title (used for type classification)
  \item \texttt{agent\_id}: which agent surface completed it
  \item \texttt{claimed\_at}: Unix timestamp when task was claimed
  \item \texttt{completed\_at}: Unix timestamp of completion
  \item \texttt{commit\_hash}: git commit closing the task (when present)
\end{itemize}

Duration $d = \texttt{completed\_at} - \texttt{claimed\_at}$ in seconds.

\todo{Measure: total N completions in done-log as of measurement date.
      Date range. Agent count.}

\subsection{Git History}

Git commits linked via \texttt{commit\_hash} provide:
file count, LOC delta, file path patterns (proxy for domain),
and co-authorship (Claude Code vs. Cursor vs. Gemini surface).

\todo{Measure: N commits with board task linkage. Coverage rate
(done-log entries with matching commit hash).}

% ─────────────────────────────────────────────────────────────────────────────
\section{Task Type Classification}
\label{sec:classification}

We classify board task titles into types using keyword matching
followed by a small LLM classifier:

\begin{table}[h]
\centering
\small
\begin{tabular}{lll}
\toprule
Type & Keywords & Example \\
\midrule
\texttt{research}    & research, synthesize, investigate, analyze & ``/research LeCun JEPA'' \\
\texttt{build}       & implement, add, build, create, scaffold   & ``Add PillarRegistry trait'' \\
\texttt{fix}         & fix, debug, resolve, repair               & ``Fix EmbeddingTrait types'' \\
\texttt{cite}        & cite, bib, reference, add citation        & ``Add RecursiveMAS to Paper 1'' \\
\texttt{paper}       & paper, tex, section, \S, draft            & ``Add JEPA section to Paper 8'' \\
\texttt{refactor}    & refactor, extract, rename, reorganize     & ``Parser W1-T2 module extraction'' \\
\texttt{review}      & review, audit, check, verify              & ``Audit paper-audit-matrix'' \\
\texttt{config}      & config, hook, script, deploy              & ``Wire SliceEmitter to GRPO'' \\
\bottomrule
\end{tabular}
\caption{Task type classification schema. \todo{Inter-rater agreement score on N=50 manual labels.}}
\end{table}

% ─────────────────────────────────────────────────────────────────────────────
\section{Duration Distribution Model}
\label{sec:system}

For each task type $\tau$, we fit a log-normal distribution to observed
durations (log-normal is standard for software task completion times~\cite{mende2010}):

\[
  \log d \sim \mathcal{N}(\mu_\tau, \sigma_\tau^2)
\]

Prediction for a new task of type $\tau$:
\begin{align}
  \hat{d} &= e^{\mu_\tau + \sigma_\tau^2/2} \quad \text{(mean estimate)} \\
  \text{CI}_{90} &= [e^{\mu_\tau - 1.645\sigma_\tau},\; e^{\mu_\tau + 1.645\sigma_\tau}]
\end{align}

For tasks with file-scope signal from the title (e.g.\ ``Add X to Papers 1, 4, 12''),
we add a file-count covariate:
$\mu_\tau \leftarrow \mu_\tau + \alpha \log(\texttt{file\_count})$
where $\alpha$ is fit from completions with git-linked file counts.

\todo{Fit parameters: $\mu_\tau$, $\sigma_\tau$ per type. Table of fitted values.}

% ─────────────────────────────────────────────────────────────────────────────
\section{The \texttt{holo\_estimate\_task\_duration} Tool}
\label{sec:tool}

\agenttime{} is deployed as an MCP tool in the \holoscript{} server.
Agents call it at board claim time with the task title and optional
file list.
The tool returns: estimated duration (seconds), 90\% CI, task type
classification, and N historical completions the estimate is based on.

\begin{lstlisting}[caption={Example tool call and response.}]
// Input
{ "task_title": "Add JEPA section to paper-8-unified-siggraph.tex",
  "files": ["research/paper-8-unified-siggraph.tex"] }

// Output
{ "type": "paper",
  "estimate_seconds": TODO,
  "ci_90": [TODO, TODO],
  "n_historical": TODO,
  "confidence": "medium" }
\end{lstlisting}

Integration: the \holomesh{} board claim hook calls
\texttt{holo\_estimate\_task\_duration} and attaches the estimate to
the claimed task. Agents use the CI to decide whether to split large
tasks before claiming.

% ─────────────────────────────────────────────────────────────────────────────
\section{Findings: Distribution Shapes by Task Type}
\label{sec:findings}

\todo{Figure: violin plot of duration distributions per task type.
      Show heavy tail for research/paper, narrow for cite/config.}

Key qualitative findings (to be quantified):
\begin{itemize}[topsep=2pt,itemsep=1pt]
  \item \textbf{Research} tasks: heavy-tailed. Long sessions are not
        outliers — they reflect genuine exploration. Clipping at the median
        would discard the most informative data points.
  \item \textbf{Cite} and \textbf{config} tasks: narrow distributions.
        Predictable within $\pm$30\% of median.
  \item \textbf{Build} tasks: bimodal. Fast (mechanical scaffolding) and
        slow (novel architecture) complete differently; single-mode models
        perform poorly.
  \item \textbf{Paper} tasks: moderate variance, correlated with file count
        (more papers to edit = longer session).
  \item \textbf{Session position effect}: first task of session is
        \todo{X\%} longer than same-type tasks after warm-up.
\end{itemize}

% ─────────────────────────────────────────────────────────────────────────────
\section{Evaluation}
\label{sec:eval}

\textbf{Setup}: hold out last 20\% of done-log entries chronologically.
Train on first 80\%. Evaluate on held-out completions.

\textbf{Metrics}:
\begin{itemize}[topsep=2pt,itemsep=1pt]
  \item MAPE (Mean Absolute Percentage Error)
  \item CI coverage: fraction of actuals falling inside 90\% CI
\end{itemize}

\textbf{Baselines}:
\begin{itemize}[topsep=2pt,itemsep=1pt]
  \item Naive mean: predict global mean duration for all tasks
  \item Type mean: predict per-type mean (no distributional fit)
  \item Story points: N/A (no human estimates in the dataset)
\end{itemize}

\todo{Results table: MAPE and CI coverage for AgentTime vs. baselines.}

% ─────────────────────────────────────────────────────────────────────────────
\section{Related Work}

\paragraph{Software task estimation.}
Story points~\cite{cohn2005}, COCOMO~\cite{boehm1981}, and
planning poker assume human cognitive load.
\agenttime{} is the first system trained on AI agent production logs.

\paragraph{Mining software repositories.}
MSR work on developer productivity~\cite{meyer2019} measures human
commit patterns. Our done-log corpus is structurally similar but
captures AI agent work cycles, which have different temporal patterns
(no circadian rhythm, no meeting interruptions, discrete context windows).

\paragraph{AI agent benchmarks.}
SWE-bench~\cite{jimenez2024} measures task success rate, not duration.
\agenttime{} is complementary: given a claimable task, predict how
long it will take.

% ─────────────────────────────────────────────────────────────────────────────
\section{Conclusion}

AI agent task duration distributions are structured and learnable.
\agenttime{} mines production done-log and git history to fit per-type
log-normal models, deployed as the \texttt{holo\_estimate\_task\_duration}
MCP tool.
The distribution shapes (heavy-tailed research, narrow cite/config, bimodal
build) are themselves a characterization of how AI agent work differs from
human software development.
\todo{Quantitative conclusion pending measurement pass.}

\bibliographystyle{ACM-Reference-Format}
\bibliography{holoscript}

\begin{thebibliography}{9}
\bibitem{cohn2005} M. Cohn. {\em Agile Estimating and Planning}. Prentice Hall, 2005.
\bibitem{boehm1981} B. Boehm. {\em Software Engineering Economics}. Prentice Hall, 1981.
\bibitem{mende2010} T. Mende, R. Koschke. Revisiting the evaluation of defect prediction models. {\em MSR}, 2010.
\bibitem{meyer2019} A. Meyer et al. Today was a good day: The daily life of software developers. {\em IEEE TSE}, 2019.
\bibitem{jimenez2024} C. Jimenez et al. SWE-bench: Can language models resolve real-world GitHub issues? {\em ICLR}, 2024.
\end{thebibliography}

\end{document}
